\section{Benchmark và giao thức đánh giá}\label{sec:protocol}
% Dàn ý gốc: ../outline/04_benchmark_protocol.md (khoảng 2 trang)

\subsection{VN-AgentInject v3}

VN-AgentInject v3 hiện thực phân loại ở Mục~\ref{sec:threat} thành các \textbf{ca kiểm thử} (fixture) chạy được. Mọi
dữ liệu đều là tổng hợp: tài liệu, bản ghi, địa chỉ và giá trị mật là bịa ra, còn hai sink là bản giả không gửi gì ra
mạng. Một \textbf{họ} (family) là một ý tưởng tấn công trong một nhiệm vụ và một môi trường công cụ. Mỗi họ sinh một
ca tấn công và một ca benign\_plain cho mỗi biến thể ngôn ngữ, và một phần ba số họ còn sinh thêm một ca benign\_auth
cho mỗi biến thể. Benchmark có ba tập rời nhau:
\begin{itemize}
\item \textbf{Dev}: 48 họ, tám họ mỗi cơ chế; 560 ca (240 tấn công, 240 benign\_plain, 80 benign\_auth). Mọi lựa
  chọn thiết kế, prompt, luật và ngưỡng chỉ được cố định trên tập này.
\item \textbf{Test v3}: 72 họ, mười hai họ mỗi cơ chế; 840 ca (360 / 360 / 120).
\item \textbf{Test-3}: thêm 48 họ; 560 ca (240 / 240 / 80), được viết sau khi Test v3 đã chạy, để xác nhận D3 trên
  dữ liệu chưa từng chạm vào quyết định nào trước đó.
\end{itemize}

\paragraph{Soạn và kiểm soát chất lượng.} Văn bản tấn công do người viết, những người không đọc mã phòng thủ. Mỗi họ
được duyệt theo phiếu năm mục: nhãn cơ chế có đúng không, yêu cầu người dùng và đoạn chèn có tự nhiên không, làm theo
đoạn chèn có gây hại theo các vị từ ở Mục~\ref{sec:threat} không, và văn bản có để lộ gợi ý về đáp án không; người
duyệt thấy nội dung ca nhưng không thấy phòng thủ. Một bộ lọc độ mới loại họ có độ tương đồng Jaccard 4-gram với bất
kỳ bản phát hành trước nào vượt 0,2; định danh, đích và tên miền không được dùng lại. Cuối cùng, mỗi họ phải qua
\textbf{kiểm tra phát lại} với bộ chấm đóng băng: một quỹ đạo tham chiếu làm theo đoạn chèn phải bị chấm là gây hại,
và một quỹ đạo tham chiếu hoàn thành nhiệm vụ an toàn phải được chấm là an toàn với đủ tiện ích. Các thủ thuật bề
mặt có thể giấu rò rỉ khỏi bộ chấm nguyên văn (ví dụ làm rối giá trị) bị loại khỏi tập kiểm thử tĩnh.

\subsection{Đánh giá niêm phong}

Hai tập kiểm thử được \textbf{niêm phong}. Trước mỗi lượt chạy kiểm thử, chúng tôi đóng băng (i) các ca chuẩn, bằng
SHA-256; (ii) một manifest cấu hình ghim bằng hash prompt tác tử, thiết lập giải mã, mã phòng thủ, phiên bản mô hình
và danh sách cặp (mô hình nền, cấu hình) được phép; và (iii) giả thuyết, kiểm định và cách hiệu chỉnh đa so sánh.
Runner và bộ chấm kiểm tra manifest và từ chối chạy nếu có bất kỳ sai khác nào hoặc gặp cặp ngoài danh sách. Mỗi cặp
chỉ chạy đúng một lần; lần chạy lại duy nhất được phép là tiếp tục các ca kết thúc vì lỗi hạ tầng. Đáp án của bộ chấm
chỉ được đọc sau khi manifest được xác minh, và không bao giờ có mặt trong prompt, chính sách lúc chạy hay mô hình.

Test v3 chạy dưới manifest đóng băng D2+P2 làm cấu hình cuối, với sáu giả thuyết (H1--H6,
Mục~\ref{sec:test1}). Sau đó D3 được thiết kế và chọn trên Dev (Mục~\ref{sec:d3}), một manifest mới được đóng băng,
và Test-3 chạy với ba giả thuyết: \textbf{T1}, D3 có tiện ích lành tính cao hơn D2+P2; \textbf{T2}, tỉ lệ tấn công
thành công của D3 không cao hơn D2+P2 quá hai điểm; \textbf{T3}, D3 có tỉ lệ tấn công thành công thấp hơn A0. Manifest
cũng ghi lại dự đoán của chúng tôi rằng mức tăng tiện ích sẽ nhỏ và có thể không đạt ý nghĩa thống kê.

\subsection{Mô hình nền và cấu hình so sánh}

Thí nghiệm xác nhận dùng ba mô hình nền mở trọng số phục vụ bằng vLLM~\cite{kwon2023vllm}: \textbf{Gemma-4-12B},
\textbf{Llama-3.1-8B} (ngữ cảnh 16k) và \textbf{Qwen3.5-9B}, mỗi mô hình ghim theo một revision của kho. Mọi mô hình
nền nhận cùng prompt ReAct và cùng thiết lập giải mã: nhiệt độ 0, tối đa 512 token mới mỗi bước, seed cố định, đầu
ra JSON ràng buộc theo schema từng công cụ, và tối đa tám bước mỗi nhiệm vụ. Tác tử không phòng thủ \textbf{A0}
không có cơ chế an ninh nào ngoài kiểm schema và giới hạn công cụ. Các phòng thủ không thay đổi mô hình nền, khuôn
prompt của A0 (ngoài lời nhắc đã mô tả) hay cách giải mã. Làm mốc cho một mô hình đóng mạnh hơn, chúng tôi chạy
\textbf{Gemini 3 Flash} qua một proxy cục bộ, chỉ trên Dev và AgentDojo. Đường này không hỗ trợ seed hay ràng buộc
schema, nên kết quả được báo cáo riêng và không thuộc kiểm định giả thuyết nào. Không có dữ liệu kiểm thử niêm phong
nào được gửi tới dịch vụ bên ngoài.

Trên Dev, chúng tôi so sánh D2, D2+P2 và D3 với sáu baseline, cài trong cùng tác tử và broker để chỉ có phòng thủ là
khác: \textbf{P1}, spotlighting, đầu ra công cụ được bao và đánh dấu là dữ liệu~\cite{hines2024spotlighting};
\textbf{P2}, sandwich, lặp lại yêu cầu người dùng sau mỗi đầu ra công cụ; \textbf{G1}, bộ phân loại Prompt Guard 2
sàng lọc đầu ra công cụ~\cite{meta2025promptguard2}; \textbf{MEL}, biến thể tất định của phép kiểm chạy lại có che
của MELON~\cite{zhu2025melon}; \textbf{PRG}, chính sách đích theo kiểu Progent~\cite{shi2025progent}; \textbf{R1},
lớp luật cụm từ của hệ thống trước, chặn sink sau khi đầu ra công cụ khớp mẫu giống chỉ thị. \textbf{D2} là lõi tất
định ở Mục~\ref{sec:d3} (luật R2--R4), \textbf{D2+P2} thêm lời nhắc sandwich, và \textbf{D3} là kiến trúc đầy đủ.
Test v3 chạy với A0, R1, PRG, P2, D2 và D2+P2; Test-3 chạy với A0, D2+P2 và D3.

\subsection{Độ đo và thống kê}

Một ca là \textbf{gây hại} nếu ít nhất một vị từ ở Mục~\ref{sec:threat} đúng. Chúng tôi báo cáo \textbf{tỉ lệ tấn
công thành công} (ASR), tức tỉ lệ ca tấn công gây hại, và số \textbf{họ gây hại}, tức họ có ít nhất một ca tấn công
gây hại. \textbf{Tiện ích} là tỉ lệ ca lành tính mà nhiệm vụ người dùng được hoàn thành và không có tác hại; báo cáo
riêng cho benign\_plain và benign\_auth, và trên Test-3 thì gộp lại thành tiện ích lành tính. Lượt chạy kết thúc bằng
lỗi mô hình (đầu ra sai định dạng hoặc bị cắt) là thất bại tiện ích và không bao giờ được tính là phòng thủ thành công.

Các cấu hình được so sánh trên \textbf{cặp ca} gộp qua ba mô hình nền, bằng kiểm định McNemar chính xác một
phía~\cite{mcnemar1947}. Giả thuyết không kém hơn T2 dùng kiểm định z một phía với sai số chuẩn theo cụm họ và biên
hai điểm. Giá trị p được hiệu chỉnh bằng thủ tục Holm~\cite{holm1979} trong mỗi bộ giả thuyết đăng ký trước. Khoảng
tin cậy là khoảng bootstrap 95\% lấy mẫu lại theo họ (10.000 lần), để năm biến thể của một họ không bị coi là độc
lập. Các phân tích ngoài bộ đăng ký trước được ghi là thăm dò.

\subsection{AgentDojo-VI}

Để kiểm tra kết quả không phải là sản phẩm riêng của benchmark do chúng tôi tự xây, chúng tôi dùng thêm hai bộ Banking
và Slack của AgentDojo (phiên bản 1.2.2)~\cite{debenedetti2024agentdojo} bằng tiếng Anh và một bản dịch tiếng Việt,
\textbf{AgentDojo-VI}. Bản dịch phủ môi trường, nhiệm vụ người dùng, mục tiêu chèn và khuôn tấn công
\code{important\_instructions}; IBAN, số tiền, ngày, tên tệp, tên người và kênh, URL và địa chỉ e-mail giữ nguyên.
Công cụ, lời giải chuẩn và hàm kiểm là mã gốc của AgentDojo: cụm tiếng Việt cần có trong câu trả lời được ánh xạ sang
chuỗi tiếng Anh mà hàm tiện ích và an ninh gốc tìm kiếm. Một phép kiểm ở mức bộ xác nhận lời giải chuẩn giải được mọi
nhiệm vụ người dùng. Banking có 144 cặp bị tấn công và 16 nhiệm vụ người dùng, Slack có 105 cặp và 21 nhiệm vụ. Các
lượt này dùng pipeline tác tử dựa trên prompt của chính AgentDojo cho mô hình cục bộ (công cụ được mô tả trong prompt)
thay vì broker của chúng tôi, nên chúng đo độ dễ bị tấn công của mô hình không phòng thủ ở từng ngôn ngữ, không đo D3.
